%!TEX program = pdflatex
% ============================================================================
%  Guhnoo Yun -- Curriculum Vitae (source of assets/pdf/Guhnoo_Yun_CV.pdf)
%
%  Build from the repo root (aux files stay out of the repo):
%    latexmk -pdf -outdir=/tmp/cvbuild cv/Guhnoo_Yun_CV.tex
%    cp /tmp/cvbuild/Guhnoo_Yun_CV.pdf assets/pdf/Guhnoo_Yun_CV.pdf
% ============================================================================
\documentclass[10pt,a4paper]{article}
\usepackage[T1]{fontenc}
\usepackage{lmodern}
\usepackage[margin=0.62in,top=0.42in,bottom=0.5in]{geometry}
\usepackage{xcolor}
\usepackage{microtype}
\usepackage{enumitem}
\usepackage{tabularx}
\usepackage{needspace}
\usepackage{hyperref}
\input{glyphtounicode}
\pdfgentounicode=1

\definecolor{cvblue}{HTML}{0E639C}
\hypersetup{colorlinks=true, urlcolor=cvblue, linkcolor=cvblue,
  pdftitle={Guhnoo Yun - Curriculum Vitae}, pdfauthor={Guhnoo Yun},
  pdfsubject={Curriculum Vitae},
  pdfkeywords={generative models, image generation, diffusion transformers, DiT, linear attention,
    spectral analysis, token mixers, computer vision, PyTorch}}

\pagestyle{empty}
\setlength{\parindent}{0pt}
\setlength{\parskip}{0pt}
\setlength{\emergencystretch}{1.5em}  % loosen a tight paragraph instead of overrunning the margin
\setlist[itemize]{leftmargin=1.8em, itemsep=0.5pt, topsep=1pt, parsep=0pt, partopsep=0pt}
\hyphenation{MMSegmentation MMDetection}

% ---- building blocks --------------------------------------------------------
\newcommand{\cvsection}[1]{\par\vspace{3pt}\needspace{5\baselineskip}%
  {\color{cvblue}\large\bfseries #1\strut\par}%
  {\color{cvblue}\hrule height 0.4pt}\vspace{2pt}}
% left text (wraps) + right-aligned bold date on the first line
\newcommand{\cvrow}[2]{\begingroup\setlength{\tabcolsep}{0pt}\noindent
  \begin{tabularx}{\linewidth}{@{}X@{\hspace{1.2em}}r@{}}#1 & \textbf{#2}\end{tabularx}\par\endgroup}
\newcommand{\cvnote}[1]{{\small\leftskip=1.6em #1\par}}
\newcommand{\cvtopic}[1]{\par\vspace{2pt}{\small\itshape #1\par}}
\newcommand{\cvsub}[1]{\par\vspace{2pt}{\small\bfseries #1\par}}
\newcommand{\hdrsep}{\unskip\hspace{1.33em}$\cdot$\hspace{1.33em}\ignorespaces}
\newcommand{\pt}[1]{{\color{cvblue}\bfseries[#1]}}                    % publication label
\newcommand{\rt}[1]{\unskip~{\color{cvblue}\scriptsize\bfseries[#1]}}  % in-text reference (tied)
\newcommand{\me}{\textbf{Guhnoo Yun}}
\newcommand{\venue}[2]{\href{#1}{\textit{#2\nocorr}}}
\newcommand{\pub}[2]{\par\vspace{3pt}\needspace{3\baselineskip}\noindent\pt{#1} #2\par}
\newcolumntype{L}{>{\raggedright\arraybackslash}X}  % ragged-right tabularx column
\newcommand{\email}[1]{\href{mailto:#1}{#1}}

\begin{document}

% ---- header -----------------------------------------------------------------
\begin{center}
  {\LARGE\bfseries Guhnoo Yun\par}\vspace{3pt}
  \href{mailto:cheeryun@gmail.com}{cheeryun@gmail.com}
  \hdrsep \href{https://scholar.google.com/citations?user=L47h01YAAAAJ}{Google Scholar}
  \hdrsep \href{https://doranlyong.github.io}{doranlyong.github.io}
  \hdrsep \href{https://github.com/DoranLyong}{GitHub}
  \hdrsep \href{https://www.linkedin.com/in/guhnoo-yun-71b40b178}{LinkedIn}\par
  Korea University \& Korea Institute of Science and Technology (KIST), Seoul, Republic of Korea\par
  \vspace{3pt}\textit{Spectral analysis and adaptive filtering for visual perception, from recognition
  to image generation.}\par
\end{center}

% ---- research focus ---------------------------------------------------------
\cvsection{Research Focus}
My research is about visual perception for machine intelligence: I analyze what vision models pass
or suppress in the spectral domain and turn that analysis into operators for visual recognition
and image generation. Next, I plan to do the same for video generation and world models at a
larger scale.
\begin{itemize}
  \item Spectral analysis of convolution and self-attention (frequency balancing, graph-spectral
        frequency responses)
  \item Spectral-adaptive operators as general-purpose replacements for softmax attention, from vision
        backbones to diffusion transformers (SPAM, SALA)
  \item Visual perception for robots and healthcare (fall detection, isolation-ward services,
        daily-life monitoring)
\end{itemize}

% ---- education --------------------------------------------------------------
\cvsection{Education}
\cvrow{\textbf{Ph.D., Computer Science and Engineering}, Korea University}{Feb 2027 (expected)}
\cvnote{Dissertation: \textit{Spectral Analysis and Adaptive Filtering for Visual Perception}}
\vspace{3pt}
\cvrow{\textbf{M.S., Mechatronics}, Gwangju Institute of Science and Technology (GIST)}{2018}
\cvnote{Thesis: \textit{Removal of Jitter Noise using a Modified Kalman Filter in 3D Shape Recovery
        from Focus}}
\vspace{3pt}
\cvrow{\textbf{B.S., Control and Measurement Engineering}, Gyeongsang National University}{2016}

% ---- research experience ----------------------------------------------------
\cvsection{Research Experience}
\cvrow{\textbf{Student Researcher}, Intelligence and Interaction Research Center, KIST}{2019--Present}
\cvnote{Supervisor: Dr.\ Dong Hwan Kim}
{\small\itshape From spectral analysis of vision models to operators for visual recognition and
image generation.\par}

\cvtopic{Linear Attention for Generative and Vision Transformers (SALAFormer; in preparation \rt{P2})}
\begin{itemize}
  \item Designed \textbf{SALA} (Spectral-Adaptive Linear Attention) as closed-form associative memory:
        each head solves a small ridge regression for what to store in its fixed-size memory,
        instead of summing raw key--value products, and a learned spectral filter on the key Gram
        matrix allocates that memory, at linear $O(N)$ inference cost.
  \item \textbf{Image generation:} lower FID than MHLA (ICLR~2026), a recent linear-attention design
        for diffusion transformers, at both model sizes: $59.80\,{\to}\,\mathbf{57.65}$ (DiT-S/2) and
        $37.47\,{\to}\,\mathbf{36.54}$ (DiT-B/2) vs.\ its published results (\mbox{class-conditional}
        ImageNet $256{\times}256$, no guidance).
  \item \textbf{Image classification:} the same mixer as a vision backbone (SALAFormer) reaches
        \textbf{77.2\%} / \textbf{81.8\%} ImageNet-1K top-1 at DeiT-Tiny/Small size, above the
        76.5\% / 81.6\% reported for ViT$^3$ (CVPR~2026) with the same recipe.
\end{itemize}

\cvtopic{Spectral Analysis of Convolution and Attention for Vision Backbones}
\begin{itemize}
  \item \textbf{SPANetV2} (TPAMI~2026): compared what convolution and self-attention pass or suppress
        in one graph-spectral framework, showing that window size shapes their frequency responses,
        and turned this into a mixer with multi-scale kernels and spectral re-scaling; the backbone
        outperforms state-of-the-art models on ImageNet-1K, COCO, and ADE20K. \rt{J1}
  \item \textbf{SPANet} (ICCV~2023): contrary to the prevailing idea of strengthening self-attention's
        high-pass filtering, showed that \emph{balancing} high and low frequencies in token mixers
        improves performance. \rt{C3}
\end{itemize}

\cvtopic{Applied Vision for Robotics, Healthcare \& 3D Shape Recovery}
\begin{itemize}
  \item Fall detection of older adults by a healthcare robot \rt{C5}; robot-assisted services in
        isolation wards \rt{C4}; \mbox{daily-life} arousal monitoring and brain--computer interface decoding
        with egocentric vision and EEG \rt{C1}\rt{C2}; a \mbox{view-consistent} evaluation metric for
        multi-view panoptic segmentation \rt{P1}.
  \item Shape from focus, following up the M.S.\ thesis: Gaussian-process regression for sampling
        from few images \rt{C6}.
\end{itemize}

\vspace{4pt}
\cvrow{\textbf{Intern Researcher}, Center for Intelligent Robotics, KIST}{2018--2019}
\begin{itemize}
  \item Perception for healthcare robots: detecting people with dementia going out, to prevent them
        from going missing \rt{C8}; accelerometer--depth fall detection for older adults \rt{C7},
        later extended with LSTMs \rt{C5}.
\end{itemize}

% ---- publications -----------------------------------------------------------
\needspace{6\baselineskip}% keep heading + first entry together
\cvsection{Publications}
\cvsub{Published / Accepted}
\pub{J1}{\me; J.~Yoo; K.~Kim; J.~Lee; P.H.~Seo; D.H.~Kim. ``Spectral-Adaptive Modulation Networks
  for Visual Perception.'' \venue{https://doi.org/10.1109/TPAMI.2026.3690455}{IEEE Trans.\ Pattern
  Anal.\ Mach.\ Intell.\ (TPAMI)}, 48(10), 11511--11525, \textbf{2026}.
  {\small\href{https://arxiv.org/abs/2503.23947}{[arXiv]}~\href{https://github.com/DoranLyong/SPANetV2-official}{[code]}}}
\pub{C1}{W.~Hyung; M.~Kim; Y.-S.~Kim; J.~Lee; B.~Ko; \me; D.H.~Kim; C.-H.~Im. ``Context-Aware
  Daily-Life Arousal Monitoring Using Egocentric Vision and EEG.''
  \venue{https://cmsworkshops.com/EMBC2026/view_paper.php?PaperNum=5021&SessionID=1198}{Annu.\ Int.\ Conf.\
  IEEE Eng.\ Med.\ Biol.\ Soc.\ (EMBC)}, \textbf{2026}.}
\pub{C2}{Y.-S.~Kim; M.~Kim; W.~Hyung; B.~Ko; \me; D.H.~Kim; C.-H.~Im. ``PeBCI: Vision-Guided
  \mbox{Brain-Computer} Interface for Context-Aware Neural Decoding in Daily Life.''
  \venue{https://cmsworkshops.com/EMBC2026/view_paper.php?PaperNum=4958&SessionID=1198}{Annu.\ Int.\ Conf.\
  IEEE Eng.\ Med.\ Biol.\ Soc.\ (EMBC)}, \textbf{2026}.}
\pub{C3}{\me; J.~Yoo; K.~Kim; J.~Lee; D.H.~Kim. ``SPANet: Frequency-balancing Token Mixer using
  Spectral Pooling Aggregation Modulation.'' \venue{https://doi.org/10.1109/ICCV51070.2023.00562}{IEEE/CVF
  Int.\ Conf.\ on Computer Vision (ICCV)}, 6090--6101, \textbf{2023}.
  {\small\href{https://openaccess.thecvf.com/content/ICCV2023/html/Yun_SPANet_Frequency-balancing_Token_Mixer_using_Spectral_Pooling_Aggregation_Modulation_ICCV_2023_paper.html}{[CVF]}
  \href{https://arxiv.org/abs/2308.11568}{[arXiv]}
  \href{https://doranlyong.github.io/projects/spanet/}{[project]}
  \href{https://github.com/DoranLyong/SPANet-official}{[code]}}}
\pub{C4}{Y.~Kwon; S.~Shin; K.~Yang; S.~Park; S.~Shin; H.~Jeon; K.~Kim; \me; \ldots; Y.~Lim. ``Heterogeneous
  Robot-Assisted Services in Isolation Wards: A System Development and Usability Study.''
  \venue{https://doi.org/10.1109/IROS55552.2023.10341857}{IEEE/RSJ Int.\ Conf.\ on
  Intelligent Robots and Systems (IROS)}, 8069--8076, \textbf{2023}.}
\pub{C5}{K.~Kim; \me; S.-K.~Park; D.H.~Kim. ``Efficient Fall Detection for a Healthcare Robot System
  Based on 3-Axis Accelerometer and Depth Sensor Fusion with LSTM Networks.''
  \venue{https://doi.org/10.1109/ICPR56361.2022.9956418}{Int.\ Conf.\ on Pattern
  Recognition (ICPR)}, 2207--2212, \textbf{2022}.}
\pub{C6}{H.-S.~Jang; \me; M.T.~Mahmood; M.-K.~Kang. ``Optimal Sampling for Shape from Focus by Using
  Gaussian Process Regression.'' \venue{https://doi.org/10.1109/ICCE46568.2020.9043150}{IEEE
  Int.\ Conf.\ on Consumer Electronics (ICCE)}, 1--4, \textbf{2020}.}
\pub{C7}{K.~Kim; \me; S.-K.~Park; D.H.~Kim. ``Fall Detection for the Elderly Based on 3-Axis
  Accelerometer and Depth Sensor Fusion with Random Forest Classifier.''
  \venue{https://doi.org/10.1109/EMBC.2019.8856698}{Annu.\ Int.\ Conf.\ IEEE Eng.\ Med.\ Biol.\ Soc.\ (EMBC)}, 4611--4614, \textbf{2019}.}
\pub{C8}{\me; K.~Kim; S.-K.~Park; D.H.~Kim. ``A Monitoring System to Support Home Health Care for the
  Elderly with Dementia by Detecting Going Out Activities Based on RGB-D Sensors.''
  \venue{https://doi.org/10.1109/URAI.2019.8768598}{Int.\ Conf.\ on Ubiquitous Robots
  (UR)}, 71--76, \textbf{2019}.}

\vspace{4pt}
\cvsub{Under Review}
\pub{P1}{Y.~Lee; B.~Ko; \me; D.H.~Kim. ``Every View Counts: View-Consistent Panoptic Quality for
  Multi-view Panoptic Segmentation.'' \venue{https://arxiv.org/abs/2610.05911}{arXiv preprint
  arXiv:2610.05911}, \textbf{2026}.}

\vspace{4pt}
\cvsub{In Preparation}
\pub{P2}{\me\ et al. ``SALAFormer: Spectral-Adaptive Linear Attention as Closed-Form Associative
  Memory'' (working title).}

% ---- awards -----------------------------------------------------------------
\cvsection{Awards \& Scholarships}
\cvrow{Outstanding Student Researcher Scholarship, KIST}{2026}
\cvrow{Outstanding Technology Award, KIST}{2023}
\cvrow{Presidential Award, 2020 International Robot Contest---Humanoid Robot Sports (Deep Learning
  Algorithm), Ministry of the Interior and Safety}{2020}

% ---- invited talks ----------------------------------------------------------
\cvsection{Invited Talks}
\cvrow{``Recent Research Trends in Computer Vision,'' seminar at Korea Polytechnic University (now
  Tech University of Korea)}{2021}
\cvrow{``6D Object Pose Estimation,'' seminar series at CUBOX (now SECERN AI)}{2021}

% ---- service ----------------------------------------------------------------
\cvsection{Professional Service}
\cvrow{Reviewer, International Conference on Pattern Recognition (ICPR)}{2024, 2026}

% ---- skills -----------------------------------------------------------------
\cvsection{Technical Skills}
\textbf{Languages \& Frameworks:} Python, PyTorch, C++, MATLAB; timm, MMDetection, MMSegmentation,
  ROS\par

% ---- references -------------------------------------------------------------
\cvsection{References}
% rows = name / affiliation / e-mail, so the e-mails stay aligned when an affiliation wraps
\noindent\begin{tabularx}{\linewidth}{@{}L@{\hspace{1em}}L@{\hspace{1em}}L@{}}
  \textbf{Dr.\ Dong Hwan Kim} & \textbf{Prof.\ Soohwan Kim} & \textbf{Dr.\ Yoonseob Lim}\\
  KIST & Kwangwoon University & KIST\\
  \email{gregorykim@kist.re.kr} & \email{kimsoohwan@kw.ac.kr} & \email{yslim@kist.re.kr}
\end{tabularx}\par

\end{document}
